% LaTeX table fragment. Requires \usepackage[utf8]{inputenc}, \usepackage{booktabs,longtable,array}, and \usepackage{pdflscape}.
\begingroup
\scriptsize
\setlength{\tabcolsep}{3pt}
\begin{landscape}
\begin{longtable}{>{\raggedright\arraybackslash}p{4.8cm} r r r r}
\caption{Phase B: H2: UNSC-Mitgliedschaft und Rohstoffabhängigkeit — Modellvergleich (1992–2023)}\label{tab:phase_b_h2_modellvergleich}\\
\toprule
Kennzahl & Basis-FE & Vollmodell-FE & Länder- und Jahres-FE & DSV-GLS
 \\
\midrule
\endfirsthead
\toprule
Kennzahl & Basis-FE & Vollmodell-FE & Länder- und Jahres-FE & DSV-GLS
 \\
\midrule
\endhead
\midrule
\multicolumn{5}{r}{Fortsetzung auf der naechsten Seite}\\
\endfoot
\bottomrule
\endlastfoot
Fokuseffekt & unsc3:resource\_dep & unsc3:resource\_dep & unsc3:resource\_dep & unsc3:resource\_dep \\
Koeffizient & 0.060 & -0.055 & -0.111 & -0.014 \\
Robuster Standardfehler & 0.078 & 0.132 & 0.103 & 0.067 \\
p-Wert & 0.439 & 0.677 & 0.286 & 0.841 \\
Beobachtungen & 380 & 212 & 212 & 212 \\
Länder & 104 & 62 & 62 & 51 \\
Länder mit unsc3 = 1 & 27 & 16 & 16 & 16 \\
Jahre & 1992--2023 & 1992--2023 & 1992--2023 & 1992--2023 \\
Länder-FE & Ja & Ja & Ja & Ja \\
Jahres-FE & Nein & Nein & Ja & Nein \\
Kontrollen & Laufzeit & DSV-Vollsatz & DSV-Vollsatz & DSV-Vollsatz \\
\end{longtable}
\noindent\footnotesize Koeffizienten der Hypothese, nicht die Koeffizienten aller Kontrollen. Heteroskedastizitaetskonsistente robuste Standardfehler. Signifikanz: * p<0.10, ** p<0.05, *** p<0.01.\par
\end{landscape}
\endgroup

